%%%%%%%%%%%%%%%%%%%%%%%%%%%%%%%%%%%%%%%%%%%%%%%%%%%%%%%%%%%%%%%%%%%%%%%%%%%%%%%%
% LateX
%
% Author: FOSSEE 
% Copyright (c) 2009, FOSSEE, IIT Bombay
%%%%%%%%%%%%%%%%%%%%%%%%%%%%%%%%%%%%%%%%%%%%%%%%%%%%%%%%%%%%%%%%%%%%%%%%%%%%%%%%

\documentclass[17pt,compress]{beamer}
\usepackage{beamerthemesplit}
\mode<presentation>
{
  \usetheme{Warsaw}
  \useoutertheme{infolines}
  \setbeamercovered{transparent}
  \setbeamertemplate{navigation symbols}{}
}
% Taken from Fernando's slides.
\usepackage{ae,aecompl}
\usepackage[scaled=.95]{helvet}

\usepackage[english]{babel}
\usepackage[latin1]{inputenc}
\usepackage[T1]{fontenc}

% change the alerted colour to LimeGreen
\definecolor{LimeGreen}{RGB}{50,205,50}
\setbeamercolor{structure}{fg=LimeGreen}
\author[FOSSEE]{}
\institute[IIT Bombay]{}
\date[]{}
% \setbeamercovered{transparent}

% theme split
\usepackage{verbatim}
\newenvironment{colorverbatim}[1][]%
{%
\color{blue}
\verbatim
}%
{%
\endverbatim
}%

\usepackage{mathpazo,courier,euler}
\usepackage{listings}
\lstset{language=sh,
    basicstyle=\ttfamily\bfseries,
  showstringspaces=false,
  keywordstyle=\color{black}\bfseries}

% logo
\logo{\includegraphics[height=1.30 cm]{../images/3t-logo.pdf}}
\logo{\includegraphics[height=1.30 cm]{../images/fossee-logo.pdf}

\hspace{7.5cm}
\includegraphics[scale=0.99]{../images/fossee-logo.pdf}\\
\hspace{281pt}
\includegraphics[scale=0.80]{../images/3t-logo.pdf}}


\begin{document}

\sffamily \bfseries
\title
[\LaTeX \: Installation]
{\LaTeX \:Installation}
\author
[FOSSEE]
{\small Talk to a Teacher\\{\color{blue}\url{http://spoken-tutorial.org}}\\\vspace{0.25cm}National Mission on Education
 through ICT\\{\color{blue}\url{ http://sakshat.ac.in}} \\ [1.65cm]
   Contributed by FOSSEE Team \\IIT Bombay  \\[0.3cm]
}

% slide 1
\begin{frame}
   \titlepage
\end{frame}

\begin{frame}
  \frametitle{Objectives}
  At the end of this tutorial you will,
  \begin{itemize}
  \item Learn to install LaTeX
  \item Learn to install a TeX editor plug-in
  \item Learn to configure the TeX editor for LaTeX
  \item Know few LaTeX packages.
  \item Be able to compile a TeX file to pdf.

  \end{itemize}
\end{frame}

\begin{frame}
  \frametitle{Prerequisites}
  Prerequisites for Installation of LaTeX
  \begin{itemize}
  \item You'll need a TeX distribution.
  \item Good text editor, DVI or PDF viewer.
  \item Updated Linux distribution as OS.
  \item Working Internet connection (for installation over the network).
  \end{itemize}
\end{frame}

\begin{frame}[fragile]
  \frametitle{TeX Distribution}
  TeX Live
  \begin{itemize}
  \item Easy way to get up and running with the TeX document production system.
  \end{itemize}
  Acquire TeX Live in many ways
  \begin{itemize}
  \item Getting TeX Live on DVD (for installation from media).
  \item Installing TeX Live over the network.
  \end{itemize}
\end{frame}

\begin{frame}[fragile]
  \frametitle{Installation}
  Running the installer
  \begin{itemize}
  \item Run the \emph{install-tl} script to install.
  \end{itemize}
  Commands
  \begin{itemize}
  \item \emph{sudo apt-get install texlive} : Basic subset of TeX Live's functionality.
  \item \emph{sudo apt-get install texlive-full} : Complete TeX Live distribution.
  \end{itemize}
\end{frame}

\begin{frame}[fragile]
  \frametitle{Installation..}
  Using a Package Manager
  \begin{itemize}
  \item Open \emph{Package Manager}
  \item Search for \emph{texlive-full}
  \item Mark for installation and apply
  \end{itemize}
\end{frame}

\begin{frame}[fragile]
  \frametitle{LaTeX Plug-in}
  Gedit has a plug-in for LaTeX which converts Gedit into a LaTeX editor
  \begin{itemize}
  \item \emph{sudo apt-get install gedit-latex-plugin} : Installs plug-in
  \item \emph Activate the plug-in : Click (Edit > Preferences > Plugins > Check LaTeX Plugin)
  \end{itemize}
\end{frame}

\begin{frame}[fragile]
  \frametitle{LaTeX Packages}
  Recommended LaTeX Packages
  \begin{itemize}
  \item \emph{latex-beamer} : to create presentations.
  \item \emph{TeXPower} : for dynamic online presentations.
  \item \emph{Prosper} : class for writing transparencies.
  \end{itemize}
\end{frame}

\begin{frame}[fragile]
  \frametitle{LaTeX Packages..}
  \begin{itemize}
  \item \emph{texlive-pictures} : for drawing graphics.
  \item \emph{texlive-latex-extra} : add-on packages
  \end{itemize}
\end{frame}

\begin{frame}[fragile]
  \frametitle{Compilation}
  LaTeX distribution installed, you may compile a LaTeX document to pdf file
  \begin{itemize}
  \item \emph{pdflatex filename.tex} : Compiles the LaTeX document.
  \item Output : PDF file.
  \end{itemize} 
\end{frame}


\begin{frame}[fragile]
	\frametitle{Summary}
	\begin{itemize}
        \item Install LaTeX on your computer.
        \item Install/Configure TeX editor with LaTeX plug-in.
        \item Choose a LaTeX package as per requirement.
	\item Compile a TeX file to PDF.
	\end{itemize}
\end{frame}

\begin{frame}[fragile]
	\frametitle{Exercise}
	\begin{itemize}
        \item How can we check the version of the LaTeX package installed ?
        \item How can we check if the plug-in is properly configured with TeX 
              editor ?
	\end{itemize}
\end{frame}

\begin{frame}[fragile]
	\frametitle{Solutions}
	\begin{itemize}
        \item Use the command \textit {latex -v} to check the version installed
        \item Edit > Preferences > Plugins > Check LaTeX Plug-in
	\end{itemize}
\end{frame}

\begin{frame}
\frametitle{SDES \& FOSSEE}
\begin{center}
\begin{itemize}
\item \small{SDES}\\
\small{\color{LimeGreen}Software Development techniques for Engineers and Scientists} \\
\scriptsize An initiative by FOSSEE. \\
\vspace{3pt}
\scriptsize For more information on SDES, please visit {\color{blue}\url{http://fossee.in/sdes}}\\
\vspace{12pt}
\item \small{FOSSEE}\\
\small {\color{LimeGreen}Free and Open-source Software for \\Science and Engineering Education} \\
\scriptsize Based at IIT Bombay, Funded by MHRD.\\
\vspace{3pt}
\scriptsize Part of National Mission on Education through ICT (NME-ICT). \\
\end{itemize}
\end{center}
\end{frame}

\begin{frame}
\frametitle{About the Spoken Tutorial Project}
\begin{itemize}
\item Watch the video available at {\color{blue}\url{http://spoken-tutorial.org /What\_is\_a\_Spoken\_Tutorial}} 
\item It summarises the Spoken Tutorial project 
\item If you do not have good bandwidth, you can download and watch it
\end{itemize}
\end{frame}

\begin{frame}
\frametitle{Spoken Tutorial Workshops}The Spoken Tutorial Project Team 
\begin{itemize}
\item Conducts workshops using spoken tutorials 
\item Gives certificates to those who pass an online test 
\item For more details, please write to \\ \hspace {0.5cm}{\color{blue}contact@spoken-tutorial.org}
\end{itemize}
\end{frame}

\begin{frame}
\frametitle{Acknowledgements}
\begin{itemize}
\item Spoken Tutorial Project is a part of the Talk to a Teacher  project 
\item It is supported by the National Mission on Education through  ICT, MHRD, Government of India 
\item More information on this Mission is available at: \\{\color{blue}\url{http://spoken-tutorial.org/NMEICT-Intro}}
\end{itemize}
\end{frame}

\begin{frame}
  \begin{block}{}
  \begin{center}
  {\Large THANK YOU!} 
  \end{center}
  \end{block}
\begin{block}{}
  \begin{center}
    For more Information, visit our website\\
    {\color{blue}\url{http://fossee.in/}}
  \end{center}  
  \end{block}
\end{frame}

\end{document}
